\corrigesChapitre{3}{Limites, continuité et équations}
\begin{corrige}{1}{Limites usuelles}
\begin{questions}
\item $+\infty$ à droite et $-\infty$ à gauche.
\item $+\infty$, car $x+1\to+\infty$ et la racine augmente sans borne.
\item $3$ : $\mathrm e^{-x}\to0$.
\end{questions}
\end{corrige}
\begin{corrige}{2}{Asymptotes}
\begin{questions}
\item $\R\setminus\{-3\}$.
\item $x=-3$ est verticale ; $y=4$ est horizontale aux deux infinis.
\item Non : le dénominateur est nul. Une asymptote verticale n’est pas un point de la courbe.
\end{questions}
\end{corrige}
\begin{corrige}{3}{Factoriser}
\begin{questions}
\item $(5+2/x)/(1+7/x)\to5$.
\item En divisant par $x^3$, le quotient tend vers $0$.
\item Pour $x\ne2$, le quotient vaut $x+2$ ; la limite est $4$.
\end{questions}
\end{corrige}
\begin{corrige}{4}{Carré et réciproque}
\begin{questions}
\item $(-1)^2=1^2$ : deux antécédents distincts ont la même image.
\item $y\mapsto\sqrt y$ sur $[0,+\infty[$ ; les graphes sont symétriques par rapport à $y=x$.
\item Le domaine est $x\geq0$ ; élever au carré donne $x=9$, qui convient.
\end{questions}
\end{corrige}
\begin{corrige}{5}{Une racine unique}
\begin{questions}
\item $f$ est continue, $f(0)=-2$ et $f(1)=1$ : le théorème des valeurs intermédiaires s’applique.
\item $f\,\!^{\prime}(x)=3x^2+2>0$ sur $\R$.
\item $f(0{,}77)=-0{,}003467<0$ et $f(0{,}78)=0{,}034552>0$ : $0{,}77<\alpha<0{,}78$.
\end{questions}
\end{corrige}
\begin{corrige}{6}{Compter les solutions}
\begin{questions}
\item Une sur chacun des deux intervalles monotones, donc deux.
\item Une sur $[-2,0]$ ; aucune sur $[0,3]$, où les valeurs sont au moins $1$.
\item Une seule : $x=0$. Le point commun aux deux intervalles ne doit pas être compté deux fois.
\end{questions}
\end{corrige}
\begin{corrige}{7}{Une inéquation avec un zéro}
\begin{questions}
\item $g\,\!^{\prime}(x)=\mathrm e^x+1>0$.
\item $g$ est continue, sa limite en $-\infty$ vaut $-\infty$, et celle en $+\infty$ vaut $+\infty$ : existence et unicité.
\item Par croissance stricte, l’ensemble des solutions est $]-\infty,\alpha]$.
\end{questions}
\end{corrige}
\begin{corrige}{8}{Précision de la dichotomie}
\begin{questions}
\item $2^{-10}=1/1024\simeq0{,}0009766$.
\item Il faut $2^n\geq10^4$ : $2^{13}=8192$ et $2^{14}=16384$, donc $14$ divisions.
\item La racine est dans l’intervalle : l’erreur sur le milieu est au plus la demi-largeur, soit $1/32768\simeq3{,}052\times10^{-5}$.
\end{questions}
\end{corrige}
\begin{corrige}{9}{La continuité est nécessaire}
\begin{questions}
\item $f(-1)=-1<0<1=f(1)$.
\item Non : la fonction ne prend que les valeurs $-1$ et $1$.
\item La continuité en $0$ manque. Un changement de signe aux bornes ne suffit pas sans continuité.
\end{questions}
\end{corrige}
\begin{corrige}{10}{Une limite trompeuse}
\begin{questions}
\item Les deux termes tendent vers $+\infty$ : forme $\infty-\infty$.
\item $h(x)=x/(\sqrt{x^2+x}+x)=1/(\sqrt{1+1/x}+1)$.
\item La limite vaut $1/(1+1)=1/2$.
\end{questions}
\end{corrige}
\begin{corrige}{11}{Newton sur un exemple}
\begin{questions}
\item $x-f(x)/f\,\!^{\prime}(x)=x-(x^2-2)/(2x)=(x+2/x)/2$ pour $x\ne0$.
\item $x_1=17/12\simeq1{,}416667$, $x_2=577/408\simeq1{,}414216$.
\item La dérivée $f\,\!^{\prime}(0)=0$ annule le dénominateur de la formule.
\end{questions}
\end{corrige}
\begin{corrige}{12}{Une recette de résolution}
\begin{questions}
\item $f(x)=\ln x+x-2$ est continue, de dérivée $1/x+1>0$ ; ses limites aux bornes sont $-\infty$ et $+\infty$.
\item $f(1{,}55)\simeq-0{,}011745$ ; $f(1{,}56)\simeq0{,}004686$. Le changement de signe et la croissance donnent l’encadrement.
\item Appliquer la fonction de dichotomie du cours sur $[1{,}55,1{,}56]$ avec précision $10^{-5}$ ; vérifier le changement de signe initial et conserver le sous-intervalle correspondant à chaque étape.
\end{questions}
\end{corrige}
